\chapter[Menius und Flacius]{Menius und Flacius im Streite über Amt und Priesterthum}
\quellseite{15}

\begin{center}
(Aus der Erlanger Zeitschrift.)
\end{center}

\begin{center}
(Schluß.)
\end{center}

Da Flacius in früheren Schriften schon sein öffentliches Auftreten mit Hinweisung auf dieselben gerechtfertigt
hatte, so sieht sich Menius veranlaßt, dieselben alsbald von seinem Gesichtspunkte aus zu erläutern. Er sagt:\footnote{Verantwortung Bog. H. Bl. 3.}

„Daß aber Illyricus seinen Beruf aus der Taufe, aus den zehn Geboten und aus dem, daß er an der Universität zu
Wittemberg eine Lectur gehabt, beweisen will, ist überaus lächerlich. — — Er sagt, er hab’ Christo in der
Tauf’ geschworen, daß er ihm wider den Teufel und allen seinen Anhang dienen, die göttliche Wahrheit bekennen
und den Satan mit aller seiner Pracht und Finanzerei verfluchen wolle.“

„Die frage ich alle verständige Christen, ja auch den Lästerer Illyricum selbst, was das heiße und sei, Christo
dienen? was das heiße und sei, die göttliche Wahrheit bekennen? was das heiße, den Satan mit seiner Pracht und
Finanzerei verfluchen? — Kann man auch Christo dienen ohne und wider sein Wort, Befehl und Ordnung? Oder ist das
nicht des Herrn Christi Wort, Befehl und Ordnung, was seine Apostel gelehret, geordnet und befohlen haben? St. Paulus
befiehlt seinem Jünger Tito, er soll in der Insel Creta die Kirchen also bestellen, daß eine jede Stadt ihren
eigenen und besonderen Bischof habe. Und Act. 20 vermahnet St. Paulus die Aeltesten zu Milet und Ephesus, sie
\quellseite{16}
sollen Acht haben beides auf sich selbst und auf die ganze Heerde, unter welche sie der Heilige Geist gesetzt habe
zu Bischöfen.“

„Da siehest du klar, wie es nach der göttlichen Ordnung, die der Herr Christus durch seine Apostel aufgerichtet
und eingesetzt hat, mit dem Beruf und Bestellung des Kirchenregiments soll gehalten werden, nämlich daß keiner
soll unberufen, ungehört und ungeprüft zu solchem Amt gelassen werden, das ist eines. Zum andern hörest du,
daß eine jede Stadt und ein jedes Kirchspiel soll seine eigene besondere Diener haben. — —“

„Will also Illyricus nach Gottes Wort, Befehl und Ordnung, wie er in der Taufe geschworen, gehandelt haben und
noch handeln, so gebühret ihm, daß er beweise, wer ihn berufen und wer ihm befohlen habe, sich des Kirchenregiments
anzunehmen und insonderheit, sich über andere Kirchendiener zu erheben, die zu rechtfertigen und zu reformiren. Item,
daß er beweise, wer ihn verhört und geprüft und tüchtig dazu erkannt habe.“

„So Illyricus von wegen seiner Taufe und der zehn Gebot pflichtig wäre, dermaßen, wie er thut, zu handeln, so
müßte unwidersprechlich folgen, daß alle getaufte Christen dergleichen auch thun müßten. Wenn nun ein jeder
Unberufene in allen Kirchen, über alle Diener richten und regieren, dieselben seines Gefallens recht sprechen
oder verdammen wollte, lieber, was wollte doch daraus werden?“

„Er gibt für, es sei aus Gottes Wort genugsam beweiset, daß alle Christen Priester sein und die Schrift auslegen
können, daß aber einer oder mehrere zum Amt erwählet werden, solches geschehe, Unordnung zu vermeiden. Damit
zeiget er an seinen hohen Verstand in christlicher Lehre, damit er in diesem Stücke beinahe so meisterlich umzugehen
weiß, als der Esel mit der Harfe. Daß alle Christen Priester sind, das ist wahr, also, daß sie an allen Orten,
zu allen Zeiten geistliche Opfer thun mögen mit Beten, Danksagen und allerlei guten Werken, item mit Geduld in
allerlei Trübsalen etc. Daß sie aber alle die Schrift auslegen können, oder allsammt die Macht haben, sich
Lehrens und Regierens in der Kirche zu unterstehen, das ist nicht. Denn die Schrift auslegen ist eine sonderliche
Gabe des Heil. Geistes, die Gott nicht Allen gibt, sondern denen allein, denen er’s geben will, 1 Cor. 14. Denn
wenn alle Christen diese Gabe hätten, was dürfte man dann des Predigtamts, welches fürnehmlich darum eingesetzt
ist, daß diejenigen, so die Gabe, die Schrift auszulegen, haben, den Andern, die solche Gabe nicht haben, dienen
und sie aus der Schrift lehren sollen? Also ist es auch nicht ein frei Ding mit den Diensten und Aemtern in der
Kirche, daß ein Jeder derselbigen sich unterstehe und die zu verwalten habe, ob sich gleich einer dünken läßt,
er könne und wolle es besser machen, denn ein Anderer, sondern es ist Gottes ernstes Gebot, Befehl und Ordnung,
daß Niemand sich einiges Amts von ihm selbst aus eigenem Thurst und Frevel unterstehen soll, er werde denn dazu
nach göttlicher Ordnung berufen und werde ihm befohlen, davon der Wende nicht viel Geschrei macht, sondern sagt
allein, daß einer oder mehrere zu Aemtern erwählet
\quellseite{17}
werden, daß geschehe, Unordnung zu vermeiden, bricht damit kurz ab, gleich als stünde es in der Menschen Wahl,
solche Ordnung zu halten oder nicht zu halten.“

Fassen wir abermals das Resultat aus der bisherigen Darlegung des Menius zusammen, so ergeben sich folgende Sätze:

1) Das Priesterthum des gläubigen Christen vollzieht sich durch geistliche Opfer mit Beten, Danksagen und allerlei
guten Werken und gibt an und für sich noch nicht Macht, die Schrift auszulegen, aus der Schrift zu lehren.

2) Die Schrift auszulegen ist eine sonderliche Gabe des Heiligen Geistes.

3) Keiner darf diese sonderliche Gabe des Heiligen Geistes ausüben, wenn er nicht dazu nach göttlicher Ordnung
berufen und ihm dieses Amt befohlen wird.

4) Das Predigtamt ist also nicht ein Amt, welches im öffentlichen Auftrage die Functionen des allgemeinen
Priesterthums öffentlich vollzieht, welches „die Aemter, so allen Christen gemein sind, von unser aller wegen
ausrichtet“\footnote{Worte Luther’s.}, sondern beruht einerseits auf einer sonderlichen Gabe des Heiligen Geistes, andererseits auf
einer göttlich geistlichen Institution, welche allein Recht und Pflicht gibt, öffentliches Zeugniß zu üben.

Menius führt zuerst zur Erhärtung seiner Ansicht eine Stelle aus Luther’s Commentar zum Galaterbrief an. Die
wesentlichen Sätze sind folgende\footnote{Siehe Walch VIII 1572 ff.}:

„Denn jetzt zu unsern Zeiten beruft uns unser Herr Gott zum Predigtamt durch Mittel, als nämlich durch Menschen.“

„Denn wiewohl ich meines Doctorates halben wohl Macht haben sollt, durchs ganze Pabstthum zu predigen, wenn man
mir’s nicht mit Gewalt wehrte, so ziemet mir’s aber doch in keinem Wege nicht, daß ich unberufen aus diesem meinem
befohlenen Kirchspiel in eine andere Stadt als ein Prediger laufen und daselbst mich Predigens unterstehen wollt.“

„Darum soll keiner dem andern in seine Ernte oder Schnitt fallen, wie der Teufel durch seine Rotten pflegt.“

„Wo aber ein Fürst oder andere Obrigkeit mich beruft oder fordert, so kann ich mit gutem und sicherem Gewissen
rühmen, daß ich je aus Gottes Befehl durch eines Menschen Stimm und Wort berufen bin.“

Eine andere Stelle Luther’s ist aus der Auslegung zum 82. Psalm, 4. B., entnommen. Da heißt es:

„Es hilft sie auch nicht, daß sie fürgeben, alle Christen sind Priester. Es ist wahr, alle Christen sind
Priester, aber nicht alle Pfarrherrn; denn über das, daß er ein Christ und Priester ist, muß er auch ein Amt
und befohlen Kirchspiel haben. Der Beruf und Befehl macht Pfarrherrn und Prediger.“

So viel ist klar: Menius beweiset mit den angeführten Stellen nichts weiteres, als daß Luther alles eigenwillige
Eingreifen und Lehren in geordneten Gemeinden als Anmaßung und wider die gottgewollte Ordnung laufend
\quellseite{18}
bezeichnet; er fordert für die Ausübung des Zeugenamtes Vollmacht und Auftrag.

Aber die Hauptfrage, auf die es ankommt und welche Menius in seiner Weise beantwortet, bleibt in den von Menius
angeführten Stellen Luther’s völlig unerledigt, nämlich: Sind die Functionen des allgemeinen Priesterthums
wesentlich verschieden von denen des Predigtamtes, besteht jenes wirklich nur darinnen, daß es die geistlichen
Opfer des Gebetes, des Dankes, der guten Werke darzubringen hat, wie Menius behauptet; oder vollzieht der Inhaber
des Predigtamtes in seinem Amte nur Functionen seines allgemeinen Christenpriesterthums?

Von der Antwort auf diese Frage bleibt immer die Stellung des Einzelnen der Gemeinde, dem Amte gegenüber, und
nebenbei auch das Recht des Flacius in Bezug auf sein öffentliches Hervortreten abhängig.

Will man sich hierüber bei Luther Antwort erholen, so sieht man sich von diesen Stellen hinweg auf eine zahllose
Menge in den übrigen Schriften Luther’s angewiesen.

Wir haben zunächst zu sehen, wie Flacius das Verhältnis des allgemeinen Priesterthums zum Amte und die Aeußerungen
Luther’s darüber auffaßt.

Vorerst leuchtet aus seiner Antwort\footnote{Apologia Bgg. C. Bl. 2 ff.} hervor, daß er keineswegs ein Verächter der kirchlichen Ordnung ist,
sondern in ihr vielmehr Gottes Ordnung und Befehl erkennt. Er rechtfertigt ein allenfallsiges Auftreten nicht
Berufener nur mit dem Falle dringender Noth.

„Was aber belanget die Vocation, ist zu wissen, erstlich daß, obwohl die ordentlichen Personen aus Gottes
Befehl sollen und müssen denen vorstehen, welchen sie verordnet sind, auch sich Niemand in ihr Amt mengen,
sondern ihnen gehorsamen und folgen soll, so ist doch solches nicht von der äußersten Noth zu verstehen. Denn
Noth, wie man sagt, bricht Eisen. Noth bricht auch Gesetz; und, wie die Juristen sagen, necessitas non habet legem.“

„Als zum Exempel, wenn einer einen erschlagen will und dergleichen oder wenn irgend ein Feuer plötzlich aufgehet,
und die ordentliche Obrigkeit entweder nicht vorhanden ist, oder sonst nicht wehren will oder auch nicht kann,
so ist allda Jedermann schuldig hinzulaufen und zu wehren, er sei gleich eine oberste oder gar unterste Person.“

„Ich frage: Da Petrus zu Antiochia also hinke und mit seinem Exempel die wahre Religion in Gefahr brachte, ob auch
ein geringster Christ, wenn er gleich ein erkaufter, leibeigener Knecht gewesen, hätte die Macht und Beruf gehabt,
Petrum öffentlich zu strafen, sonderlich so es Paulus und andere hohe Lehrer nicht gethan hätten? Freilich hätte
er’s nicht allein Macht gehabt, sondern wäre es schuldig gewesen. Ob auch solcher leibeigene Knecht damit sich
über Petrum erhoben und eine Unordnung oder Rottenei angerichtet hätte? ich halte nicht.“
\quellseite{19}
„Auf daß wir aber näher zur Sache kommen: Es ist eine löbliche und christliche Gewohnheit, daß in der Noth
auch die Weiber mögen taufen. Item, es hat Dr. Martinus seliger Gedächtniß „in Captivitate Babylonica“,
im Büchlein „von dreien Mauern des Papsts“, item „über die Epistel St. Petri“ und im Büchlein „de
instituendis Ministris“ reichlich bewiesen, daß alle Christen rechte Priester sind, daß sie auch mögen die
Schlüssel gebrauchen, alle Lehre urtheilen c. Doch solches alles von der großen hohen Noth zu verstehen ist,
und daß keine Unordnung und Zerstörung in der Kirche geschehe; denn es muß ja alles zur Erbauung dienen.“

„Ich hab nichts mehr gethan, denn daß ich in solcher großen gefährlichen Noth mit meiner großen Gefahr die
Wahrheit schriftlich bekannt und aus der Schrift bewiesen und wiederum die Falschheit gestraft und widerlegt,
auch die Leute bei der Wahrheit zu verharren und die Lügen und Irrthümer zu meiden vermahnet habe.“

„Solches haben vorzeiten wie auch zu unsern Zeiten sehr viel Ungeweihete oder Unordinirte gethan c.“

„Ich weiß nicht anders, wiewohl ich nicht streiten will, denn daß auch Philippus am ersten allein zu der
griechischen Lection berufen worden ist\footnote{Melanchthon wurde erst am 19. Sept. 1519 zum Baccalaureus der Theologie ernannt und als solcher in die theologische Facultät aufgenommen. Aber schon während des vorhergehenden Jahres las er über den Brief an Titus. G. Matthes, Ph. Melanchthon, S. 35.}, wie er denn auch darum von den Wittenbergischen Bürgern und Bauern der
Grek genannt worden ist; gleichwohl hat er auch Theologiam zu lesen angefangen, ja auch zu ordiniren die Prediger,
so er doch nicht zum Kirchendienst ordiniret ist worden, welches gleichwohl in andern Kirchen nicht gewöhnlich.“

„Es ist jetzt gut Menio zu schreyen und zu schreiben von der Vocation; aber zur Zeit des Interims und der
Adiaphoristerei war schier Niemand so sehr geweibet und geschmieret, so geschickt oder so hochgelehrt, der gern
hätte wollen die Wahrheit Gottes klar bekennen, die mancherlei Interim verdammen und widerlegen, geduldiglich
dafür leiden c.“

So ist es also die Noth, welche Anlaß sein kann, sich der Lehre und Vermahnung aus Gottes Wort anzunehmen: aber
inneren Antrieb, Recht und Gewalt in solchem Falle gibt die Zugehörigkeit zu Christus, die priesterliche Würde
der Christen.

„Erstlich entsagen alle Christen in der Taufe dem Teufel und allen seinen Werken und dagegen ergeben und verpflichten
sie sich Christo, daß sie ihm allein dienen, seine Ehre und des Nächsten Heil, aufs beste sie immer können,
suchen und für die Ehre Christi wider den Teufel und alle seine Gewalt aufs heftigste fechten und streiten wollen.“

„Ja“, sagt Menius, „die Taufe bringt nicht mit sich, daß du dich erheben solltest über alle Kirchen und
Prediger und sie deines Gefallens regieren und reformiren.“
\quellseite{20}
„Antwort. Ich erhebe mich über Niemand. Ich gebiete Niemand etwas, ich verbiete auch nichts. Das ist nicht
sich über alle Kirchen und Schulen erheben, wenn gleich Paulus Petrum strafet, oder auch ein gemeiner Christ den
andern. Matth. 18.“

„Zum andern, so hab ich aus den zehn Geboten in Summa ein hartes Gebot und ernsten Befehl, daß ich Gott soll
lieben von ganzem Herzen und meinen Nächsten als mich selbst. Ich muß lieben nicht allein mit Worten, sondern auch
mit der That, also daß ich zur Zeit der Noth mein Leben für ihn lasse. Womit hab’ ich nun in solchen bösen
Zeiten und äußerster Noth besser können beide Gott und meinen Nächsten lieben, denn eben in der Verantwortung
seiner himmlischen heilsamen Wahrheit und Bestrafung, auch Widerlegung der mancherlei geschwinden Irrthümer?“

„Ich hab’ nicht lang können und sollen mit mir disputiren, wer doch mein Nächster sei. Denn ich hab’
ihn vor mir gehabt, und gesehen so viel tausend pusillos Christi oder arme Gewissen, von den Interimisten und
Adiaphoristen greulich nicht allein verwundet, sondern auch getödtet.“

„Zum dritten, sagt Paulus: daß Gottes Wort unter uns Christen reichlich soll wohnen und daß wir uns einer den
andern vermahnen und trösten sollen mit der Schrift Col. 3. Da ist ja ein klarer Befehl Gottes, daß ein Christ
wohl Macht habe, ja auch schuldig sei, den andern zu stärken und zu lehren, auch die Verführer zu strafen.“

„Zum vierten: Christus Matth. 18. befiehlt Jedermann insonderheit die Schlüssel, da er saget: So dein Bruder
sündiget \&c. Vide Luther. in libro de instituendis ministris super hunc locum.“

„Zum fünften: Es schreibet auch Paulus 1 Cor. 14.: daß wir alle prophetiren d.\,i. lehren können, und befehlt
auch, daß, so dem Sitzenden etwas offenbar würde, soll der Andre schweigen. Solches versteht Dr. Martinus Luther
im Büchlein de instituendis ministris von allen Christen. Wie denn auch Paulus solche Epistel nicht allein an
alle Geweihte, sondern an alle Corinther geschrieben hat. Derwegen so hab’ ich auch Macht als ein Christ.“

„Zum sechsten: Es ist in den zehn Geboten eins, das heißt also: Du sollst nicht falsch Zeugniß geben. Welches
ebenso viel ist und heißt als: Du sollst der Wahrheit Zeugniß geben. — — Wir haben ein ernst Gebot, daß
wir Abgötterei und falsche Lehrer fliehen sollen. Wie können wir aber das thun, wenn wir nicht richten sollen
alle Lehrer und Lehre? — —“

„Zum siebenten: Wir haben auch ein sehr strenges Gebot, daß wir Christum und seine Lehre bekennen sollen und
müssen. Welches so gar nöthig ist, daß Junker Major darf schreiben: wir werden dadurch selig. Die Bekenntniß aber
fasset in sich beide, das Lob der Wahrheit und die Verdammung der Falschheit. Was ist nun das für ein Bekenntniß,
wenns nicht öffentlich geschiehet? — —“

„Zum achten: Es gebeut auch Gott durch Salomo Prov. 24: Wenn
\quellseite{21}
wir sehen, daß einer unbillig zum Tode geführt wird, daß wir ihm allerdings helfen, ihn retten und unsere Hand
nicht abziehen sollen. So man nun so gar sehr schuldig ist, diejenigen, so zum leiblichen Tode geführet werden,
zu retten, wie viel mehr diejenigen, so zum geistlichen, ewigen Tode der Seelen durch falsche Lehrer gerissen
werden? — —

„Zum ersten: Es zeigten die Apostel Christo an Marc. 9. und Luc. 9.: daß einer da wäre, welcher die Teufel im
Namen Christi austriebe, so er doch Christo nicht folgete, und baten, er wollte ihm’s verbieten. Aber Christus
wollt es ihm nicht allein nicht verbieten, sondern bestätiget sein Thun und saget: Wer nicht wider uns ist, der
ist für uns. Item wer in meinem Namen die Teufel austreibet, der kann nicht bald übel von mir reden. Damit denn
Christus anzeigt, daß ihm wohlgefalle, daß die Leute, sie seien auch wer sie wollen, nur mit dem geringsten Dienst
sein Reich mehren und ausbreiten helfen. Also schreibet auch Paulus Phil. 1.: daß etliche Christen predigen aus
gutem Herzen, etliche aus bösem, etliche auch unrecht, Paulo zum Verdrieß. Er aber sei froh, daß nur Christus
gepredigt wird.“ — —

Im Uebrigen führt nun Flacius noch aus, wie er doch auch nicht ohne allen menschlichen oder durch Menschen
geschehenen Beruf zum Streite wider die Irrthümer hervorgetreten sei: Er sei ein Studiosus der heiligen Schrift,
habe zu Basel, Tübingen und Wittenberg Theologie studirt; er sei ein Magister gewesen und habe als solcher geloben
müssen, nicht allein Philosophium, sondern auch die wahre Religion treulich lehren und fördern zu wollen; er sei
ein Lehrer der heil. Schrift des alten Testamentes in seiner natürlichen Sprach zu Wittenberg gewesen; endlich
sei er durch viele angesehene Männer, vor allen durch den trefflichen Superintendenten von Hamburg Dr. Aepinus
in diesem seinem Thun bestärkt worden.

Zum Schlusse wendet sich Flacius noch gegen die falsche Deutung, welche Menius den oben angeführten Stellen aus
Luthers Schriften gibt:

„Daß aber Menius die einfältigen Christen gern betrügen wollte mit etlichen Sprüchen Dr. Martini vom Beruf,
weil er nichts Klares aus Gottes Wort aufbringen kann, ist zu wissen, daß solche Sprüche alle dahin gehen,
daß ein Unberufener sich nicht soll in ein fremdes Amt eindringen, auch sonst keine Unordnung mit seinem Thun
anrichten außerhalb der äußersten Noth, welches ich nicht gethan habe. Darum so gehen mich solche Sprüche
nicht an. Ich hab’ aber andere viele und klare Sprüche Lutheri, da er klar wider die Papisten beweiset, daß
alle Christen Priester sein, und nur wohl Macht haben zu lehren, aber gleichwohl also, daß sie Unordnung und
Aergerniß außerhalb der äußersten Noth verhüten.“

„Man lese sein Buch de instituendis Ministris, an die Böhmen geschrieben, darinnen er klar bekennet und saget,
daß er in vielen Schriften wider die Papisten bewiesen habe, daß wir Alle Priester sein, Allen die Schlüssel
überantwortet sein, Alle mögen lehren und predigen. Dergleichen im Buch de captivitate Babylonica, im Buch von
dreien Mauern des Papstes, in der
\quellseite{22}
Schrift von guten Werken, in der Auslegung der ersten Epistel Petri, da er klar saget, daß alle Christen
Priester sein, alle Macht haben zu predigen; aber daß gleichwohl etliche das Amt verwesen, geschehe der Ordnung
und Geschicklichkeit halben. Dergleichen im Buch vom Mißbrauch der Messe handelt er diese Materie klarlich,
daß alle Christen mögen predigen und lehren. Dergleichen Zeugnisse sind fast alle Bücher Dr. Martini voll.“

„Aus welchen Zeugnissen Lutheri ist je klar, daß alle Christenmenschen Priester sein, wohl Macht und Beruf
von Gott haben, das Wort Gottes zu lehren, sonderlich die da tüchtig sind; daß aber etliche zum Amt erwählet
werden, ist die Ursach, daß nicht alle tüchtig sind, auch nicht alle zu lehren stets Zeit und Zuhörer haben,
und auf daß eine nützliche Ordnung in der Kirche Gottes sei und erhalten werde.“ So weit Flacius.

\begin{unterschrift}
W. Präger.
\end{unterschrift}
